%! Modeling with Exponential Functions: Interest, Half-Life, and e — Unit 6, Lesson 6.4 — Group Activity
\documentclass[10pt]{article}
\usepackage{algebra2-article}
\usepackage{algebra2-boxes}
\newcommand{\TallMath}[1]{$\displaystyle #1\rule[-1.4em]{0pt}{3.2em}$}

\begin{document}

\pageheader{Unit 6, Lesson 6.4}{Group Activity}
\namepartnerperiod

\vspace{0.08in}

\begin{headlinebox}{forestbg}
\small\textbf{The Base Is One Period. The Exponent Counts Periods.} Every model below is
$y=ab^{x}$ wearing a costume. Before either partner touches a calculator, both of you must agree out
loud on two things: \textbf{what is one period here?} and \textbf{how many of them fit in the time
given?} Answer those and the substitution writes itself. Round money to the nearest cent and
milligrams to two decimal places.
\end{headlinebox}

\vspace{0.06in}

\begin{tcolorbox}[colback=white, colframe=black!40, title=\textbf{Tier R --- Remediate},
    fonttitle=\bfseries, arc=1.5mm, left=3mm, right=3mm, top=2mm, bottom=2mm, breakable]
\textbf{Part 1 --- Set up before you solve.} \$2000 is invested at $3\%$ for $10$ years. Fill in one
row at a time; the last column is the whole point.
\begin{center}
\small
\renewcommand{\arraystretch}{1.6}
\begin{tabular}{|l|c|c|c|c|}
\hline
\textbf{Compounded} & $n$ & \textbf{base} $1+\frac{r}{n}$ & \textbf{exponent} $nt$ & \textbf{$A=2000\left(1+\frac{r}{n}\right)^{nt}$} \\
\hline
annually  & $1$ & $1.03$ & $10$ & \blank{2.4cm} \\
\hline
quarterly & $4$ & \blank{1.6cm} & \blank{1.4cm} & \blank{2.4cm} \\
\hline
\end{tabular}
\end{center}

\vspace{2pt}
Compounding four times as often earned an extra \blank{2.0cm} over ten years. Was that more or less
than your partner expected? \blank{5.0cm}

\vspace{6pt}
\textbf{Part 2 --- A medicine that halves.} A $600$ mg dose has a half-life of $12$ hours.

\vspace{3pt}
\hspace*{1.0em} Model: $A(t)=$ \blank{4.0cm}, \; where $t$ is measured in \blank{1.8cm}.

\vspace{3pt}
\begin{center}
\small
\renewcommand{\arraystretch}{1.4}
\begin{tabular}{|l|c|c|c|c|c|}
\hline
$t$ (hours) & $0$ & $12$ & $24$ & $36$ & $48$ \\ \hline
halvings $t/12$ & $0$ & \blank{1.0cm} & \blank{1.0cm} & \blank{1.0cm} & \blank{1.0cm} \\ \hline
$A(t)$ (mg) & $600$ & \blank{1.3cm} & \blank{1.3cm} & \blank{1.3cm} & \blank{1.3cm} \\ \hline
\end{tabular}
\end{center}

\vspace{2pt}
Every row cuts the one before it in half, yet the amount never reaches $0$. Which feature of the graph
says so? \blank{6.0cm}

\vspace{5pt}
\textbf{Part 3 --- One period, named.} For each model, say what one period is and how many periods fit
in the time asked for.

\vspace{2pt}
\begin{center}
\small
\renewcommand{\arraystretch}{1.6}
\begin{tabularx}{\linewidth}{@{}l X X@{}}
\hline
\textbf{Model} & \textbf{One period is\dots} & \textbf{Periods in the time given} \\
\hline
$4000(1.02)^{4t}$, \; $t=3$ yr & \blank{3.0cm} & \blank{2.6cm} \\
$50\left(\frac12\right)^{t/9}$, \; $t=27$ days & \blank{3.0cm} & \blank{2.6cm} \\
$300\cdot2^{\,t/4}$, \; $t=20$ hours & \blank{3.0cm} & \blank{2.6cm} \\
\hline
\end{tabularx}
\end{center}
\end{tcolorbox}

\begin{tcolorbox}[colback=white, colframe=black!40, title=\textbf{Tier A --- Approaching Proficiency},
    fonttitle=\bfseries, arc=1.5mm, left=3mm, right=3mm, top=2mm, bottom=2mm, breakable]
\textbf{Part 1 --- Two offers, one deposit.} A credit union offers \$5000 at $3.6\%$ compounded
\emph{monthly}. A bank offers the same \$5000 at $3.6\%$ compounded \emph{continuously}. Both are
left alone for $8$ years.

\vspace{3pt}
\begin{center}
\small
\renewcommand{\arraystretch}{1.7}
\begin{tabularx}{\linewidth}{@{}l X X c@{}}
\hline
& \textbf{Formula used} & \textbf{Substitution} & \textbf{Balance} \\
\hline
Credit union & \blank{3.0cm} & \blank{4.0cm} & \blank{2.2cm} \\
Bank & \blank{3.0cm} & \blank{4.0cm} & \blank{2.2cm} \\
\hline
\end{tabularx}
\end{center}

\vspace{2pt}
\textbf{(a)} Which is better, and by how much? \blank{5.0cm}

\vspace{3pt}
\textbf{(b)} A classmate says continuous compounding should ``obviously'' win by a lot, since it
compounds infinitely often. Using your two numbers, explain what is actually true.
\writelines{2}

\vspace{5pt}
\textbf{Part 2 --- Between the halvings.} A $120$ mg dose has a half-life of $5$ hours.

\vspace{2pt}
\hspace*{1.0em} $A(t)=$ \blank{3.6cm} \hfill $A(10)=$ \blank{2.0cm} \hfill $A(17)=$ \blank{2.0cm}

\vspace{3pt}
$A(17)$ is not a whole number of halvings. Explain why that is fine, and check that your answer is the
size it ought to be by naming the two table values it must fall between. \writelines{2}

\vspace{5pt}
\textbf{Part 3 --- Error analysis.} An account pays $8\%$ compounded quarterly. Three students model
\$1000 after $5$ years. Name each error, then write the correct model and value.

\vspace{2pt}
\textbf{(a)} Renata writes $A=1000(1.08)^{20}$. \writelines{2}

\vspace{1pt}
\textbf{(b)} Miles writes $A=1000(1.02)^{5}$. \writelines{2}

\vspace{1pt}
\textbf{(c)} Correct model: $A=$ \blank{4.2cm} $=$ \blank{2.4cm}

\vspace{3pt}
Renata and Miles made \emph{opposite} halves of the same mistake. Say what the one sentence is that
both of them broke. \blank{7.0cm}
\end{tcolorbox}

\begin{tcolorbox}[colback=white, colframe=black!40, title=\textbf{Tier E --- Extension},
    fonttitle=\bfseries, arc=1.5mm, left=3mm, right=3mm, top=2mm, bottom=2mm, breakable]
\textbf{Part 1 --- Where the chopping stops.} You saw $\left(1+\frac1m\right)^{m}$ close in on $e$.
Now watch the same thing happen at a realistic rate: $\left(1+\frac{0.06}{n}\right)^{n}$, the growth
factor for one year at $6\%$.

\vspace{3pt}
\begin{center}
\small
\renewcommand{\arraystretch}{1.4}
\begin{tabular}{|l|c|c|c|c|c|}
\hline
$n$ & $1$ & $4$ & $12$ & $365$ & continuous \\
\hline
one-year factor & $1.060000$ & $1.061364$ & $1.061678$ & $1.061831$ & \blank{1.9cm} \\
\hline
\end{tabular}
\end{center}
\hspace*{1.0em} The last cell is $e^{\,0.06}$. Compute it to six decimal places: \blank{2.4cm}

\vspace{3pt}
\textbf{(a)} A bank advertises ``compounded every \emph{second}.'' How much more than daily
compounding could that possibly pay on \$1000 in one year? \blank{4.0cm}

\vspace{3pt}
\textbf{(b)} Why is $e^{\,0.06}$ a \emph{ceiling} rather than just the next entry in the table?
\writelines{2}

\vspace{6pt}
\textbf{Part 2 --- The Rule of 72.} Bankers estimate doubling time with a shortcut: at $r\%$ per year
compounded annually, money doubles in about $\frac{72}{r}$ years. Test it against the graph-reading you
did in the notes.

\vspace{3pt}
\begin{center}
\small
\renewcommand{\arraystretch}{1.6}
\begin{tabularx}{\linewidth}{@{}c X c X c@{}}
\hline
\textbf{Rate} & \textbf{Rule of 72 estimate} & \textbf{$(1+r)^{t}$ at $t=$ estimate} & \textbf{Over or under 2?} & \textbf{True doubling time} \\
\hline
$5\%$ & \blank{1.8cm} & \blank{2.0cm} & \blank{1.6cm} & about $14.2$ yr \\
$8\%$ & \blank{1.8cm} & \blank{2.0cm} & \blank{1.6cm} & about $9.0$ yr \\
$12\%$ & \blank{1.8cm} & \blank{2.0cm} & \blank{1.6cm} & about $6.1$ yr \\
\hline
\end{tabularx}
\end{center}

\vspace{2pt}
\textbf{(a)} At which of the three rates is the shortcut most accurate? \blank{2.6cm}

\vspace{2pt}
\textbf{(b)} A shortcut is only worth having if you know when to distrust it. Based on your table, when
does the Rule of 72 start to go wrong? \writelines{2}

\vspace{6pt}
\textbf{Part 3 --- Which of these can you actually finish?} Each question below reduces to an
exponential equation. Reduce it, then decide whether Lesson 6.3's common-base method finishes the job.

\vspace{3pt}
\begin{center}
\small
\renewcommand{\arraystretch}{1.8}
\begin{tabularx}{\linewidth}{@{}l X c X@{}}
\hline
\textbf{Question} & \textbf{Equation, after dividing} & \textbf{Common base?} & \textbf{If yes, the answer} \\
\hline
$600$ mg, $h=12$ h: when is $75$ mg left? & \blank{2.8cm} & \blank{1.2cm} & \blank{2.4cm} \\
$300$ cells, $d=4$ h: when are there $2400$? & \blank{2.8cm} & \blank{1.2cm} & \blank{2.4cm} \\
\$2000 at $3\%$ annually: when is it \$4000? & \blank{2.8cm} & \blank{1.2cm} & \blank{2.4cm} \\
$600$ mg, $h=12$ h: when is $200$ mg left? & \blank{2.8cm} & \blank{1.2cm} & \blank{2.4cm} \\
\hline
\end{tabularx}
\end{center}

\vspace{2pt}
\textbf{(a)} Two of the four defeated the method. For each one, say whether the answer \emph{exists},
and how you know without solving. \writelines{2}

\vspace{2pt}
\textbf{(b)} Rows 1 and 4 are the same drug with the same half-life. One works and one does not. What
is it about the \emph{target number} that decides it? \writelines{2}
\end{tcolorbox}

\end{document}
